% ==================== 4 单技术方案与联用决策 ====================
\section{单技术方案与联用决策}

\subsection{只能选一种时的方案}

若只保留一台仪器，方案仍应以红外为骨架，因为本问题问的是配位方式，红外对羧酸根最敏感。方案不能照搬前文，需改以下几处：KBr 压片换成 ATR 或 DRIFTS，避开吸潮与散射；1520--1560 cm$^{-1}$ 的重叠峰用差谱与分峰拆开；配位判据用纯相 TiO$_2$、油酸钠与油酸标样重新标定。低波数 Ti—O—Ti 只能给晶相半定量\cite{piao2026ir}。若只剩拉曼，则用 785 nm、低功率、长积分，先测相组成，再由 C—H/C=C 判断包覆层；而配位方式只能作倾向性判断，不能给出确切的结论。

\subsection{联用能多知道什么}

两种谱叠在同一视场上，才能表明清楚配位与晶相是否相关。红外只表明油酸怎样结合，拉曼只表明哪里是哪种相；把二者对齐后，若某一晶相的 $\Delta\nu$ 或构象不同，就可判断油酸是否优先稳定该晶相。这种“配位—晶相—位置”的对应，单一技术不能表明清楚\cite{kurouski2020}。
